\section{What repeated readings add to measurement}
\label{gap29:measurement}

A bank may ask an assistant to recover financial distress from customer
notes before assessing the effect of a servicing action. Reading the same
notes several times is attractive because disagreement is easy to observe.
Agreement can be less informative: several readers may share the same
omission. The distinction affects whether text can help control confounding.
The literature on text-based causal adjustment identifies measurement error,
collider information and post-treatment content as separate problems.
\citep{exec13_census_text_review_arxiv}

\subsection{A common error creates a lower limit}
Consider a project model, not an empirical claim about any language model.
Let latent distress be $U$, the action index $A=U+\varepsilon_A$, and
the outcome $Y=\tau A+U+\varepsilon_Y$, with $\tau=-1/5$. The latent
variable and the two disturbances are independent, centred, unit-variance
Gaussians. Suppose reading $j$ produces $W_j=U+e_j$. The measurement errors
are independent of $U,\varepsilon_A,\varepsilon_Y$, have variance $s$ and
equal pairwise correlation $\rho$. These premises make the calculation
transparent; actual extraction errors may violate all of them.

The averaged reading has error variance
\[
 s_m=\operatorname{Var}\!\left(\frac{1}{m}\sum_{j=1}^m e_j\right)
     =s\left(\rho+\frac{1-\rho}{m}\right).
\]
The covariance of $(A,\overline W)$ is
$\left(\begin{smallmatrix}2&1\\1&1+s_m\end{smallmatrix}\right)$,
and its covariance with $Y$ is $(2\tau+1,\tau+1)'$. Solving these two
normal equations gives the coefficient on $A$ after adjustment:
\[
 \beta_{A\mid\overline W}=\tau+\frac{s_m}{1+2s_m}.
\]
Under exchangeability, the contrasts between readings contain only noise
orthogonal to both $A$ and $Y$. Including every $W_j$ therefore yields the
same population coefficient as adjusting for their mean. Exact rational
calculations agree with the full $(m+1)$-variable covariance solution in all
20 prespecified cases.

For $\tau=-1/5$, the adjusted coefficient is negative only if $s_m<1/3$.
With independent unit-variance errors, at least four readings cross that
threshold. With error correlation $1/5$, at least seven are needed. If
$s\rho\geq 1/3$, no finite number can repair the sign. The measurement
variance tends to $s\rho$, so repeated readings retain a shared error floor.
At $s=1$ and $\rho=1/2$, the limiting adjusted coefficient is $1/20$,
although the true effect is $-1/5$.

\begin{figure}[htbp]
\centering
\includegraphics[width=.94\textwidth]{scientific-gap-closure/2026-10-29/measurement-floor.pdf}
\caption{Project-derived population coefficients under repeated noisy
measurements. Curves are exact calculations under the stated Gaussian model.
The true coefficient is -0.2; positive coefficients have the wrong sign.
No language model was evaluated.}
\label{gap29:measurement-floor}
\end{figure}

\subsection{An evaluation should test the proposed repair}
This calculation gives a reason to measure errors against an independent
reference rather than to count agreement between prompts. A validation sample
could compare a structured dictionary, a single extraction, repeated
extractions and retrieval-assisted extraction using the same permitted notes
and decision date. Each method should record omitted information, timing
violations and error dependence across cases. An outcome observed after the
action must not quietly enter the pretreatment measurement.

The economic test is whether the measured improvement changes the causal
conclusion at an acceptable cost. Lower text-classification error alone is
insufficient if the remaining error tracks the action or outcome. Conversely,
the shared-error limit suggests a concrete remedy: obtain an independent
measurement or better source information, rather than simply requesting more
readings of the same incomplete account. A second model is not automatically
an independent measurement. Independence requires evidence.

Graph construction raises the same issue. An assistant can propose
mechanisms and competing diagrams, but the relevant evidence must concern
the bank's assignment and information process. A formal method can evaluate
an intervention under a supplied graph; it cannot establish the graph's
missing arrows from fluent prose. The next prototype comparison should
therefore count consequential formulation and graph errors, alongside time
spent correcting them. It should keep a questionnaire plus ordinary search
as an alternative and preserve the formal method's unsupported conclusions.

The present derivation establishes a limit and a discriminating experiment.
It measures no live model benefit, does not close the unread literature
census and supplies no universal measurement-error correction. Its practical
contribution is to identify which additional information could make repeated
assistance useful, and which repetition cannot supply.
